
\section{Implementation and Training}
\label{sec:experiments}
% \red{Suggestion: Do not mix experiments, experiment setup, training details, configuration, etc., with results of the experiments and discussion.}

 To ensure efficient training, gradients for both classical and quantum parameters were calculated using PyTorch's automatic differentiation. The training was carried out with a batch size of 64 and a maximum of 20,000 epochs in all case studies except in the cavity problem, where we used 12500 maximum iterations. The Adam optimizer was employed to ensure stable convergence. These settings were uniformly applied across all study cases to provide a consistent basis for comparison.  Additionally, training data was selected through the Sobol sequence for the cavity problem and random sampling for all other cases.

While other training configurations, such as increasing the number of iterations or adjusting the batch size, could yield better test performance, the goal is to evaluate the models under identical conditions. Given the near-infinite hyperparameter space, this controlled evaluation ensures a fair assessment of the relative strengths and limitations of each approach. It is also important to acknowledge that the quantum experiments conducted in this study were performed using simulations rather than actual quantum hardware. These simulations inherently operate at significantly reduced speeds compared to their classical counterparts due to the substantial computational overhead associated with emulating quantum circuits on classical machines. Future implementations on actual quantum hardware would eliminate this simulation overhead, potentially revealing additional performance characteristics not captured in the current study.


Finally, the experimental results indicate that the computations on the CPU were faster than those on the GPU. This is likely due to our study's relatively small batch size, circuit complexity, and limited number of qubits/qumodes: five qubits for DV-based models and two qumodes for CV-based models. Moreover, the quantum backends employed in our experiments are not optimized for GPU efficiency. Given these limitations, we opted to run our experiments on a CPU system with 6148 CPU @ 2.40GHz and 192GB of RAM.

Further details on training configurations and implementation strategies are provided in the following subsections.

\subsection{CV-based QCPINN}

To implement and train our hybrid CV-based QCPINN model, we employed the PennyLane framework \cite{bergholm2018pennylane} with the Strawberry Fields Fock backend \cite{killoran2019strawberry}. Each CV-based QNN layer was configured for two qumodes with a cutoff dimension of 20. Quantum parameters were initialized using controlled random distributions, with small values for active parameters (e.g., $r_{\text{squeeze}}$, $r_{\text{disp}}$) to ensure numerical stability. Classical layers were initialized using Xavier initialization.
  
Despite the advantages of CV quantum models, their operations must be carefully managed to maintain numerical stability. Squeezing operations, for instance, are highly sensitive to large values, as excessive squeezing amplifies quadrature fluctuations, increasing photon number variance and truncation errors in Fock-space representations~\cite{lvovsky2015squeezed}. Similarly, displacement operations must be constrained to prevent quantum states from exceeding computational cutoffs, which can introduce simulation errors~\cite{ferraro2005gaussian}. Likewise, Kerr interactions introduce phase shifts that scale quadratically with photon numbers, making them highly sensitive to large values~\cite{drummond1980quantum}. Therefore, careful optimization strategies are necessary to ensure stable training. 

In our experiments, we employed a learning rate of $1.0 \times 10^{-4}$ to account for the model’s sensitivity to gradient updates. A \texttt{ReduceLROnPlateau} scheduler was implemented with a patience of 20 epochs, reducing the learning rate by a factor of 0.5 when no improvement was observed, down to a minimum of $1.0 \times 10^{-6}$. Additionally, we applied gradient clipping to prevent exploding gradients, constraining their norm to a maximum of 0.1.
  
To further stabilize the CV transformations, we scaled the training parameters by small factors. While this mitigates instability and prevents divergence in quantum state evolution, it can introduce the vanishing gradient problem. This issue is well-documented in classical neural network literature and arises when excessively small initialization values diminish gradients during backpropagation, slowing or even stalling the learning process. In our case, scaling parameters such as $r_{\text{squeeze}}$, $r_{\text{disp}}$, and ${\text{kerr}}$ to small values restrict the network's ability to propagate meaningful gradients effectively. 
%TODO: needs more research
% The situation is further complicated by the barren plateau phenomenon~\cite{sharma2022trainability,volkoff2021large}, which is particularly severe in CV quantum systems due to their infinite-dimensional Hilbert space and continuous parameterization.

\subsection{DV-based QCPINN}

The optimization method used in the DV-based QCPINN models relies on the Adam optimizer with a learning rate of 0.005. Additionally, a learning rate scheduler, \textit{ReduceLROnPlateau}, is employed to adjust the learning rate when the loss plateaus dynamically. Specifically, the scheduler reduces the learning rate by a factor of 0.9 if no improvement is observed for 1000 consecutive epochs, promoting stability and convergence. Furthermore, the training function employs gradient clipping to prevent exploding gradients, clipping them according to their norm to ensure they do not exceed a unity value. During training, Pauli-Z expectation measurements are performed on each qubit to extract relevant information about the system and guide the optimization process.




\subsection{Classical PINN Model}

The classical baseline model consists of three main components: a preprocessing network, a hidden network, and a postprocessing network. Although the pre/postprocessing networks remain consistent with the hybrid QCPINN framework, the hidden network varies between configurations.
    \begin{itemize}
        \item \textbf{Model-1}: The hidden network consists of two fully connected layers, each containing 50 neurons with a Tanh activation function between layers.
        \item \textbf{Model-2}: The hidden network is entirely omitted, reducing the model to a simpler structure.
    \end{itemize}
    
These two configurations represent different levels of model complexity, with Model-1 containing more layers and, consequently, a higher number of trainable parameters compared to Model-2. We use these classical models as benchmarks to evaluate and compare the performance of the proposed hybrid models in figure~\ref{fig:architecture}. Training settings, including optimizer and batch size, were kept consistent to enable a fair comparison between classical and quantum-enhanced approaches. The learning rate is $0.005$


Finally, we implemented five distinct PDEs, including the Helmholtz equation, the time-dependent 2D cavity problem, the 1D wave equation, the nonlinear Klein-Gordon equation, and a convection-diffusion equation. The mathematical formulation, boundary conditions, and corresponding loss function designs for each PDE are detailed in~\ref{app:appendix_b}.

    

